
\exercise
Suppose $T \in \L(V\1, W)$. Show that with respect to each choice of bases of $V$ and $W\3$, the matrix of $T$ has at least $\dim \ran T$ nonzero entries.


\exercise
Suppose $T \in \L(V\1, W)$, where $V$ and $W$ are finite-dimensional and nonzero. Prove that $\dim \ran T = 1$ if and only if there exist a basis of $V$ and a basis of $W$ such that with respect to these bases, all entries of $\M(T)$ equal $1$.


\exercise
Suppose $v_1, \ldots, v_n$ is a basis of $V$ and $w_1, \ldots, w_m$ is a basis of $W\3$.
\begin{subtheorem}
\item
Show that if $S, T \in \L(V\1, W)$, then $\M(S+T) = \M(S) + \M(T)$.
\item
Show that if $\la \in \F{}$ and $T \in \L(V\1, W)$, then $\M(\la T) = \la \M(T)$.
\end{subtheorem}
\exercisecomment{This exercise asks you to verify \ref{101} and \ref{102}.}


\exercise
Suppose that $D \in \L\bigl(\P\1_3(\R{}), \P\1_2(\R{}) \bigr)$ is the differentiation map defined by \mbox{$Dp = p'\!$}. Find a basis of $\P\1_3(\R{})$ and a basis of $\P\1_2(\R{})$ such that the matrix of $D$ with respect to these bases is\index{differentiation linear map}
\[
\mat{cccc}{
1 & 0 & 0 & 0 \\
0 & 1 & 0 & 0 \\
0 & 0 & 1 & 0 \\}.
\]
\vspace{-.15in}
\exercisecomment{Compare with Example \ref{diffmatrix}. The next exercise generalizes this exercise.}


\exercise
Suppose $V$ and $W$ are finite-dimensional and $T \in \L(V\1, W)$. Prove that there exist a basis of $V$ and a basis of $W$ such that with respect to these bases, all entries of $\M(T)$ are $0$ except that the entries in row $k$, column~$k$, equal $1$ if $1 \le k \le \dim \ran T\3$. \label{zeroonematrix}


\exercise
Suppose $v_1, \dots, v_m$ is a basis of $V$ and $W$ is finite-dimensional. Suppose $T \in \L(V\1, W)$. Prove that there exists a basis $w_1, \dots, w_n$ of $W$ such that all entries in the first column of $\M(T)$ [with respect to the bases $v_1, \dots, v_m$ and $w_1, \dots, w_n$] are $0$ except for possibly a $1$ in the first row, first column.
\exercisecomment{In this exercise, unlike Exercise \ref{zeroonematrix}, you are given the basis of $V$ instead of being able to choose a basis of $V\1$.}


\exercise
Suppose $w_1, \dots, w_n$ is a basis of $W$ and $V$ is finite-dimensional. Suppose $T \in \L(V\1, W)$. Prove that there exists a basis $v_1, \dots, v_m$ of $V$ such that all entries in the first row of $\M(T)$ [with respect to the bases $v_1, \dots, v_m$ and $w_1, \dots, w_n$] are $0$ except for possibly a $1$ in the first row, first column.
\exercisecomment{In this exercise, unlike Exercise \ref{zeroonematrix}, you are given the basis of $W$ instead of being able to choose a basis of $W\3$.}


\exercise
Suppose $A$ is an \by{m}{n} matrix and $B$ is an \by{n}{p} matrix. Prove that \label{3think1}
\[
(AB)_{j,\cdot} = A_{j,\cdot} \, B
\]
for each $1 \le j \le m$. In other words, show that row $j$ of $AB$ equals (row~$j$~of~$A$) times $B$.
\exercisecomment{This exercise gives the row version of \ref{3matrixcol}.}


\exercise
Suppose $a = \mat{ccc}{a_1 & \cdots & a_n}$ is a \by{1}{n} matrix and $B$ is an \by{n}{p} matrix. Prove that \label{3think2}
\[
aB = a_1 B_{1, \cdot} + \dots + a_n B_{n, \cdot}\ .
\]
In other words, show that $aB$ is a linear combination of the rows of $B$, with the scalars that multiply the rows coming from $a$.
\exercisecomment{This exercise gives the row version of \ref{3combcol}.}


%\exercise
%Suppose $A$ is an \by{n}{n} matrix and $m$ is an integer greater than~$2$. Show that the entry in row~$j$, column %$k$ of $A^m$ (which is defined to mean $\underbrace{A A \dotsm A}_{m \text{\ times}}$) is
%\[ \sum_{p_1=1}^n \sum_{p_2=1}^n \dotsm \sum_{p_{m-1}=1}^n A_{j, p_1} A_{p_1, p_2} A_{p_2, p_3} \dotsm A_{p_{m-1}, k}.= \]
%\exercisedisplayend

\exercise
Give an example of \by{2}{2} matrices $A$ and $B$ such that $AB \ne BA$. \label{noncommute}\index{commutativity}


\exercise
Prove that the distributive property holds for matrix addition and matrix
multiplication.  In other words, suppose $A$, $B$, $C$, $D$, $E$, and $F$ are matrices whose
sizes are such that $A(B+C)$ and $(D+E)F$ make sense.  Explain why $AB + AC$ and $DF + EF$ both make sense
and prove that\label{3distrib}\index{distributive property}
\[
A(B+C) = AB + AC \andd (D+E)F = DF + EF.
\]
\exercisedisplayend


\exercise
Prove that matrix multiplication is associative.  In other words, suppose $A$,
$B$, and $C$ are matrices whose sizes are such that $(AB)C$ makes sense.
Explain why $A(BC)$ makes sense and prove that\label{assoc} \[(AB)C = A(BC).\]
\exercisecomment{Try to find a clean proof that illustrates the following quote from Emil Artin: ``It is my experience that proofs involving matrices can be shortened by 50\% if one throws the matrices out.''\index{Artin, Emil}}


\exercise
Suppose $A$ is an \by{n}{n} matrix and $1 \le j, k \le n$. Show that the entry in row~$j$, column~$k$, of $A^3$ (which is defined to mean $AAA$) is
\[
\sum_{p=1}^n \sum_{r=1}^n A_{j, p} \, A_{p, r} \, A_{r, k}.
\]
\vspace{-10bp}


\exercise
Suppose $m$ and $n$ are positive integers. Prove that the function $A \mapsto \tran{A}$ is a linear map from $\F{m, n}$ to
$\FF{n, m}$.\label{3translinear}


\pagebreak[3]

\exercise
Prove that if $A$ is an \by{m}{n} matrix and $C$ is an \by{n}{p} matrix, then\label{3tran}
\[
\tran{(AC)} = \tran{C} \tran{A}\3.
\]
\exercisedisplayend
\exercisecomment{This exercise shows that the transpose of the product of two matrices is the product of the transposes in the opposite order.}


\pagebreak[3]

\exercise
Suppose $A$ is an \by{m}{n} matrix with $A \ne 0$. Prove that the rank of $A$ is $1$ if and only if there exist $(c_1, \dots, c_m) \in \F{m}$ and $(d_1, \dots, d_n)\in \F{n}$ such that $A_{j,k} = c_j d_k$ for every $j = 1, \dots, m$ and every $k = 1, \dots, n$.


\exercise
Suppose $T \in \L(V)$, and $u_1, \dots, u_n$ and $v_1, \dots, v_n$ are bases of $V\1$. Prove that the following are equivalent.
\begin{subtheorem}*
\item
$T$ is injective.
\item
The columns of $\M(T)$ are linearly independent in $\FF{n,1}$.
\item
The columns of $\M(T)$ span $\FF{n,1}$.
\item
The rows of $\M(T)$ span $\FF{1,n}$.
\item
The rows of $\M(T)$ are linearly independent in $\FF{1,n}$.
\end{subtheorem}
Here $\M(T)$ means $\M\bigl( T, (u_1, \dots, u_n), (v_1, \dots, v_n) \bigr)$.


